% !TeX program = lualatex
% !TeX encoding = utf8
% !TeX spellcheck = uk_UA
% !TEX indexprogram = upmendex
% !TeX root =../ClassicalElectrodynamics.tex

% chktex-file 3
% chktex-file 7
% chktex-file 12
% chktex-file 25

%=========================================================
\chapter{Реакція випромінювання}\label{\currfilebase}
%=========================================================

Досі при розгляді руху зарядженої частинки в зовнішньому полі ми нехтували впливом власного поля частинки на її рух. Такий підхід є послідовним у нерелятивістській границі, де власна енергія частинки незмінна, та й у загальному випадку є коректним наближенням, коли зміна кінетичної енергії за час порядку $\tau_0 = q^2/(mc^3)$ мала порівняно з кінетичною енергією. Проте ситуація принципово змінюється, коли частинка рухається з прискоренням: вона випромінює~\eqref{eq:total_power_nonrel}, і це випромінювання переносить енергію та імпульс від частинки до поля. Виникає питання: яка саме сила уповільнює рух і відбирає в частинки енергію, що йде на випромінювання?

%% --------------------------------------------------------
\section{Нерелятивістська сила Абрагама–Лоренца}
%% --------------------------------------------------------

Розглянемо нерелятивістський рух частинки із зарядом $q$ та масою $m$. Потужність, що витрачається на
дипольне~випромінювання:
%>> - - - - - - - - - - - - - - - - - - - - - - - - - - - - - - - - - - - - ->>
\begin{equation}\label{eq:Larmor}
	N = \frac{2q^2}{3c^3}\,\dot{\vect{v}}^2.
\end{equation}
%<< - - - - - - - - - - - - - - - - - - - - - - - - - - - - - - - - - - - - -<<
За законом збереження енергії ця потужність дорівнює роботі деякої ефективної сили $\vect{f}_\text{r}$ на одиницю
часу:
%>> - - - - - - - - - - - - - - - - - - - - - - - - - - - - - - - - - - - - ->>
\begin{equation}\label{eq:power_balance}
	\vect{f}_\text{r} \cdot \vect{v} = -N.
\end{equation}
%<< - - - - - - - - - - - - - - - - - - - - - - - - - - - - - - - - - - - - -<<
Цю силу називають \emph{силою реакції випромінювання}. Відношення~\eqref{eq:power_balance} визначає лише проєкцію $\vect{f}_\text{r}$ на вектор швидкості, але не дає її значення в довільний момент часу. Однак для руху з обмеженим прискоренням і швидкістю (наприклад, для фінітного руху або для коливань) можна вимагати виконання балансу~\eqref{eq:power_balance} у середньому за часом. Знайдемо вигляд функції \(\vect{f}_\text{r}\). Враховуючи формулу Лармора
\eqref{eq:Larmor}, можна записати необхідну умову у вигляді

\begin{equation*}
\int\limits_{t_1}^{t_2} \vect{f}_\text{r} \cdot \vect{v} \, dt = -\int\limits_{t_1}^{t_2} \frac{2q^2}{3c^3} \,
\dot{\vect{v}} \cdot \dot{\vect{v}} \, dt.
\end{equation*}

Беручи другий інтеграл частинами, отримуємо

\begin{equation*}
\int\limits_{t_1}^{t_2} \vect{f}_\text{r} \cdot \vect{v} \, dt =
\frac{2q^2}{3c^3} \bigl[\vect{v} \cdot \dot{\vect{v}}\bigr]_{t_1}^{t_2} -
\frac{2q^2}{3c^3} \int\limits_{t_1}^{t_2} (\vect{v} \cdot \ddot{\vect{v}}) \, dt.
\end{equation*}

Якщо рух періодичний або такий, що скалярний добуток \(\vect{v} \cdot \dot{\vect{v}}\) дорівнює нулю
при \(t = t_1\) та \(t = t_2\), то останнє співвідношення можна переписати у вигляді

\begin{equation*}
\int\limits_{t_1}^{t_2} \left( \vect{f}_\text{r} - \frac{2q^2}{3c^3} \ddot{\vect{v}} \right) \cdot \vect{v} \, dt = 0.
\end{equation*}

Таким чином, як силу реакції випромінювання можна прийняти величину:


%Проінтегруємо ліву частину рівняння за частинами:
%%>> - - - - - - - - - - - - - - - - - - - - - - - - - - - - - - - - - - - - ->>
%\begin{equation*}
%	\int\limits_{t_1}^{t_2} \dot{\vect{v}}^2\, dt
%	= \left[\vect{v}\cdot\dot{\vect{v}}\right]_{t_1}^{t_2}
%	- \int\limits_{t_1}^{t_2} \vect{v}\cdot\ddot{\vect{v}}\, dt.
%\end{equation*}
%%<< - - - - - - - - - - - - - - - - - - - - - - - - - - - - - - - - - - - - -<<
%Для фінітного руху перший член зникає в середньому (або дорівнює нулю, якщо інтегрувати за періодом). Тому умова~\eqref{eq:power_balance} може задовільнитись, якщо покласти:

%>> - - - - - - - - - - - - - - - - - - - - - - - - - - - - - - - - - - - - ->>
\begin{equation}\label{eq:AL_force}
	\vect{f}_\text{r} = \frac{2q^2}{3c^3}\,\ddot{\vect{v}}
	=  m\tau_0 \dot{\vect{a}},
\end{equation}
%<< - - - - - - - - - - - - - - - - - - - - - - - - - - - - - - - - - - - - -<<
де $\vect{a} = \dot{\vect{v}}$ --- прискорення, \(\tau_0 = \dfrac{2q^2}{3mc^3}\) --- характерний час. Це є \emph{формулою Абрагама–Лоренца}.
Рівняння руху частинки з урахуванням зовнішньої сили $\vect{F}$ та реакції
випромінювання набуває вигляду:
%>> - - - - - - - - - - - - - - - - - - - - - - - - - - - - - - - - - - - - ->>
\begin{equation}\label{eq:EOM_AL}
	m\dot{\vect{v}} = \vect{F} + \frac{2q^2}{3c^3}\,\ddot{\vect{v}}.
\end{equation}
%<< - - - - - - - - - - - - - - - - - - - - - - - - - - - - - - - - - - - - -<<

%Коефіцієнт при старшій похідній швидкості визначає характерний масштаб часу $\tau_0$, властивий самій зарядженій частинці:
%%>> - - - - - - - - - - - - - - - - - - - - - - - - - - - - - - - - - - - - ->>
%\begin{equation}\label{eq:tau_0}
%	\tau_0 = \frac{2q^2}{3mc^3}.
%\end{equation}
%<< - - - - - - - - - - - - - - - - - - - - - - - - - - - - - - - - - - - - -<<
Для електрона характерний час становить $\tau_0 \approx 6.3 \cdot 10^{-24}$~с. Просторовий масштаб, який з точністю до чисельного множника визначає \emph{класичний радіус електрона}:
%>> - - - - - - - - - - - - - - - - - - - - - - - - - - - - - - - - - - - - ->>
\begin{equation}\label{eq:r_0}
	r_0 = \frac{q^2}{mc^2} \approx 2.8 \cdot 10^{-13}\text{~см}.
\end{equation}


%% --------------------------------------------------------
\section{Труднощі класичної теорії}
%% --------------------------------------------------------

Рівняння~\eqref{eq:EOM_AL} є диференціальним рівнянням \emph{третього порядку} відносно координати, і тому містить
незвичну фізичну поведінку, що не має аналога в ньютонівській механіці.

\emph{Саморозгінні розв'язки.} При $\vect{F} = 0$ рівняння~\eqref{eq:EOM_AL} має розв'язок
$\vect{v}(t) = \vect{v}_0 \exp(t/\tau_0)$, тобто прискорення зростає необмежено з часом навіть за відсутності
зовнішніх сил. Такі розв'язки є нефізичними і мають відкидатися.

\emph{Передприскорення.} Формально для вибору єдиного розв'язку рівняння третього порядку потрібна не лише початкове положення та швидкість, а й початкове прискорення. Якщо вимагати відсутності саморозгінних розв'язків, частинка починає заздалегідь <<відчувати>> прикладену силу ще до того, як вона починає діяти, на час порядку $\tau_0$. Цей ефект передприскорення не можна визначити в експериментах (оскільки $\tau_0$ нескінченно мале порівняно з будь-якими лабораторними масштабами), але вказує на внутрішню суперечливість класичної теорії точкового заряду.

Розглянуті труднощі є відображенням того факту, що класична електродинаміка неспроможна самоузгоджено описати структуру власного поля точкового заряду.


%% --------------------------------------------------------
\section{Релятивістське узагальнення}
%% --------------------------------------------------------

Для отримання релятивістського аналога~\eqref{eq:AL_force} зауважимо, що шукана 4-сила $f^{\mu}$ повинна:
\begin{itemize}
	\item бути ковектором\footnote{Тобто перетворюватися як контраваріантний 4-вектор при перетвореннях Лоренца.}
	      і будуватися з 4-швидкості $u^{\mu} = dx^{\mu}/ds$ та 4-прискорення $w^{\mu} = du^{\mu}/ds$;
	\item задовольняти умові ортогональності до 4-швидкості: $f_{\mu}u^{\mu} = 0$, що випливає з~\eqref{eq:2.6};
	\item переходити у~\eqref{eq:AL_force} в нерелятивістській границі.
\end{itemize}
З цих умов єдиний можливий вираз є:
%>> - - - - - - - - - - - - - - - - - - - - - - - - - - - - - - - - - - - - ->>
\begin{equation}\label{eq:LAD_force}
	f^{\mu} = \frac{2q^2}{3c}
	\left( \frac{dw^{\mu}}{ds} + w^{\nu}w_{\nu}\, u^{\mu} \right).
\end{equation}
%<< - - - - - - - - - - - - - - - - - - - - - - - - - - - - - - - - - - - - -<<
Це \emph{сила Лоренца–Абрагама–Дірака}. Ортогональність легко перевіряється безпосередньо: $f^{\mu}u_{\mu} = \frac{2q^2}{3c}(u_{\mu}dw^{\mu}/ds + w^{\nu}w_{\nu})$, а оскільки $d(w^{\mu}u_{\mu})/ds = (dw^{\mu}/ds)\,u_{\mu} + w^{\mu}w_{\mu} = 0$, то $f^{\mu}u_{\mu} = 0$. Нерелятивістська границя відновлює формулу~\eqref{eq:AL_force}, що може бути перевірено безпосередньою підстановкою.

%% --------------------------------------------------------
\section{Наближення Ландау–Ліфшиця}
%% --------------------------------------------------------

Для практичних розрахунків за умови, що реакція випромінювання є малою
поправкою до руху під дією зовнішньої сили, зручним є наближення,
запропоноване Ландау та Ліфшицем~\cite{LandauLifshitz2003}. Параметр малості ---
відношення $\tau_0/T \ll 1$, де $T$ --- характерний час зміни зовнішнього поля.

Повне рівняння руху з урахуванням сили~\eqref{eq:LAD_force}:
%>> - - - - - - - - - - - - - - - - - - - - - - - - - - - - - - - - - - - - ->>
\begin{equation}\label{eq:full_EOM}
	mc\,\frac{du^{\mu}}{ds} = \frac{q}{c}F^{\mu\nu}u_{\nu}
	+ \frac{2q^2}{3c}\left(\frac{dw^{\mu}}{ds} + w^{\nu}w_{\nu}\,u^{\mu}\right).
\end{equation}
%<< - - - - - - - - - - - - - - - - - - - - - - - - - - - - - - - - - - - - -<<
Оскільки другий доданок малий, у ньому достатньо використати нульове
наближення для 4-прискорення, тобто покласти
%>> - - - - - - - - - - - - - - - - - - - - - - - - - - - - - - - - - - - - ->>
\begin{equation}\label{eq:w_zero}
	w^{\mu} = \frac{q}{mc^2}F^{\mu\nu}u_{\nu}.
\end{equation}
%<< - - - - - - - - - - - - - - - - - - - - - - - - - - - - - - - - - - - - -<<
Другий доданок у~\eqref{eq:LAD_force} обчислюється підстановкою~\eqref{eq:w_zero}
безпосередньо:
%>> - - - - - - - - - - - - - - - - - - - - - - - - - - - - - - - - - - - - ->>
\begin{equation}\label{eq:ww}
	w^{\nu}w_{\nu} = \frac{q^2}{m^2c^4}\,
	\bigl(F^{\nu\alpha}u_{\alpha}\bigr)\bigl(F_{\nu\beta}u^{\beta}\bigr).
\end{equation}
%<< - - - - - - - - - - - - - - - - - - - - - - - - - - - - - - - - - - - - -<<
Для першого доданка диференціюємо~\eqref{eq:w_zero} по $s$:
%>> - - - - - - - - - - - - - - - - - - - - - - - - - - - - - - - - - - - - ->>
\begin{equation}\label{eq:dw_ds}
	\frac{dw^{\mu}}{ds}
	= \frac{q}{mc^2}\frac{d}{ds}\bigl(F^{\mu\nu}u_{\nu}\bigr)
	= \frac{q}{mc^2}\left(
		\partial_{\alpha}F^{\mu\nu}\,u^{\alpha}u_{\nu}
		+ F^{\mu\nu}\frac{du_{\nu}}{ds}
	  \right).
\end{equation}
%<< - - - - - - - - - - - - - - - - - - - - - - - - - - - - - - - - - - - - -<<
У другому доданку~\eqref{eq:dw_ds} знову підставляємо нульове наближення
$du_{\nu}/ds = w_{\nu}$, визначене формулою~\eqref{eq:w_zero}:
%>> - - - - - - - - - - - - - - - - - - - - - - - - - - - - - - - - - - - - ->>
\begin{equation}\label{eq:Fw}
	F^{\mu\nu}w_{\nu}
	= \frac{q}{mc^2}F^{\mu\nu}F_{\nu\alpha}u^{\alpha}.
\end{equation}
%<< - - - - - - - - - - - - - - - - - - - - - - - - - - - - - - - - - - - - -<<
Підставляючи~\eqref{eq:ww}--\eqref{eq:Fw} у вираз для сили~\eqref{eq:LAD_force},
отримуємо \emph{наближення Ландау–Ліфшиця}:
%>> - - - - - - - - - - - - - - - - - - - - - - - - - - - - - - - - - - - - ->>
\begin{equation}\label{eq:LL_force}
	f^{\mu}_\text{ЛЛ} = \frac{2q^3}{3mc^3}
	\left[
		\partial_{\alpha}F^{\mu\nu}\,u^{\alpha}u_{\nu}
		+ \frac{q}{mc^2}F^{\mu\alpha}F_{\alpha\nu}u^{\nu}
		+ \frac{q}{m^2c^4}
		  \bigl(F_{\alpha\nu}u^{\nu}\bigr)\!\bigl(F^{\alpha\beta}u_{\beta}\bigr)
		  u^{\mu}
	\right].
\end{equation}
%<< - - - - - - - - - - - - - - - - - - - - - - - - - - - - - - - - - - - - -<<
Три доданки мають прозорий зміст: перший враховує неоднорідність поля
вздовж траєкторії, другий --- подвійну дію поля на 4-швидкість, третій
забезпечує ортогональність $f^{\mu}_\text{ЛЛ}u_{\mu}=0$, яку легко
перевірити безпосередньо. Сила~\eqref{eq:LL_force} не породжує
саморозгінних розв'язків і передприскорення, оскільки містить лише
першу похідну за власним часом, а не третю, як рівняння~\eqref{eq:EOM_AL}.
Вона є коректним описом реакції випромінювання у всіх практичних
ситуаціях, де $\tau_0 \ll T$.

%% --------------------------------------------------------
\section*{Задачі}
%% --------------------------------------------------------

\begin{problem}
	Довести, що в нерелятивістській границі 4-сила~\eqref{eq:LAD_force} зводиться до формули Абрагама–Лоренца
	\eqref{eq:AL_force}.
    %
	\emph{Вказівка.} Покласти $u^{\mu} \approx (1, \vect{v}/c)$, $w^{\mu} = du^{\mu}/ds \approx (0, \dot{\vect{v}}/c^2)$,
	враховуючи, що $ds \approx c\,dt$ при $v \ll c$.
\end{problem}

\begin{problem}
	Заряджений гармонічний осцилятор із зарядом $q$, масою $m$ та власною частотою $\omega_0$ знаходиться під впливом
	сили реакції випромінювання. Показати, що в рівнянні руху з'являється «затухання», і знайти декремент загасання.

%	\emph{Відповідь.} Ефективне рівняння руху
%	$\ddot{x} + \gamma\dot{x} + \omega_0^2 x = 0$,
%	де $\gamma = \tau_0\omega_0^2 = \dfrac{2q^2\omega_0^2}{3mc^3}$ --- декремент загасання.
\end{problem}

\begin{problem}\label{prb:runaway}
	Записати загальний розв'язок однорідного рівняння~\eqref{eq:EOM_AL} (при $\vect{F}=0$) і знайти умову на початкові
	дані, що виключає саморозгінний розв'язок.

%	\emph{Відповідь.} Загальний розв'язок $\dot{\vect{v}}(t) = \vect{A} e^{t/\tau_0} + \vect{B}$, де $\vect{A}$,
%	$\vect{B}$ --- сталі вектори. Умова відсутності саморозгону: $\vect{A} = 0$, тобто $\dot{\vect{v}}(0) = \vect{B}$ ---
%	довільне, але прискорення не зростає з часом.
\end{problem}

\begin{problem}
	Частинка із зарядом $q$ та масою $m$ рухається в однорідному магнітному полі $\Bfield$ по колу.
	Знайти час, за який частинка гальмується до зупинки внаслідок втрат на випромінювання, вважаючи нерелятивістський
	рух і початкову швидкість $v_0 \ll c$.
    %
	\emph{Вказівка.} Скористатись формулою~\eqref{eq:Larmor} та законом збереження енергії.
	Прискорення при рівномірному русі по колу: $|\dot{\vect{v}}| = \omega_c v$, $\omega_c = qB/(mc)$.

%	\emph{Відповідь.}
%	$t_{\text{гальм}} = \dfrac{3m^3c^5}{2q^4B^2}\cdot\dfrac{1}{v_0^2}\cdot\left(1 - \dfrac{v^2}{v_0^2}\right)^{-1}$
%	у початковий момент, або повна кінетична енергія $\mathcal{E}_k(t) = \dfrac{\mathcal{E}_{k0}}{1+t/t^*}$, де
%	$t^* = \dfrac{3m^3c^5}{4q^4B^2 v_0^{-2}}$.
\end{problem}

% ============================================================

\begin{problem}\label{prb:LL_crossed}
    Релятивістська заряджена частинка з масою $m$ та зарядом $q$ рухається
    в однорідних схрещених полях $\Efield \perp \Bfield$, $E < B$.
    Знайти силу реакції випромінювання в наближенні Ландау--Ліфшиця
    та показати, що вона має компоненту, перпендикулярну до швидкості,
    яка спричиняє дрейф траєкторії у напрямку $\Efield\times\Bfield$.

    \emph{Вказівка.} Для однорідних полів $\partial_{\alpha}F^{\mu}{}_{\nu}=0$.
    Скористатись лоренц-інваріантом $F_{\alpha\nu}F^{\alpha\nu} = 2(\Bfield^2 - \Efield^2) > 0$.


% ============================================================
%  РОЗВ'ЯЗОК
% ============================================================

%\vspace{1ex}

%\emph{Розв'язок}
%
%\textbf{1. Спрощення сили ЛЛ.}
%Поля однорідні, тому перший доданок у~\eqref{eq:LL_force} тотожно зникає,
%$\partial_{\alpha}F^{\mu}{}_{\nu} = 0$, і сила ЛЛ набуває вигляду:
%%
%\begin{equation}\label{eq:LL_crossed_simplified}
%    f^{\mu}_\text{ЛЛ} = \frac{2q^4}{3m^2c^5}
%    \left[
%        F^{\mu}{}_{\alpha}F^{\alpha}{}_{\nu}u^{\nu}
%        - \mathcal{I}\,u^{\mu}
%    \right],
%\end{equation}
%%
%де введено лоренц-скаляр
%%
%\begin{equation}\label{eq:I_def}
%    \mathcal{I} \equiv \bigl(F_{\alpha\nu}u^{\nu}\bigr)
%                       \bigl(F^{\alpha}{}_{\beta}u^{\beta}\bigr)
%                     = (Fu)^{\alpha}(Fu)_{\alpha}.
%\end{equation}
%%
%Зауважимо, що $\mathcal{I}$ --- це \emph{не} $F_{\alpha\nu}u^{\alpha}u^{\nu}$,
%яке тотожно нуль через антисиметрію $F_{\alpha\nu}=-F_{\nu\alpha}$,
%а квадрат 4-вектора лоренц-сили.
%
%\textbf{2. Система координат і тензор поля.}
%Покладемо $\vect{B} = B\,\hat{z}$, $\vect{E} = E\,\hat{x}$.
%У гауссовій системі ($F^{0i}=E^i$, $F^{ij}=-\varepsilon^{ijk}B_k$)
%ненульові компоненти матриці $F^{\mu}{}_{\nu}$:
%%
%\begin{equation}\label{eq:F_components}
%    F^{0}{}_{1} = E, \quad
%    F^{1}{}_{0} = -E, \quad
%    F^{1}{}_{2} = B, \quad
%    F^{2}{}_{1} = -B.
%\end{equation}
%
%\textbf{3. Лоренц-інваріанти поля.}
%Перший інваріант:
%%
%\begin{equation}\label{eq:inv1}
%    F_{\alpha\nu}F^{\alpha\nu} = 2(B^2 - E^2) \equiv 2\mathcal{F} > 0,
%\end{equation}
%%
%що підтверджує «магнітний» характер поля при $E < B$.
%Другий інваріант $\tfrac{1}{2}\varepsilon^{\alpha\nu\lambda\rho}F_{\alpha\nu}F_{\lambda\rho}
%= 4\,\vect{E}\cdot\vect{B} = 0$, оскільки $\vect{E}\perp\vect{B}$.
%
%\textbf{4. Вектор лоренц-сили $(Fu)^{\mu}$.}
%При $u^{\mu} = \gamma(1,\,v_x/c,\,v_y/c,\,0)$:
%%
%\begin{align}\label{eq:Fu}
%    (Fu)^{0} &= E\,\gamma v_x/c, \notag\\
%    (Fu)^{1} &= -E\gamma + B\gamma v_y/c, \notag\\
%    (Fu)^{2} &= -B\gamma v_x/c, \notag\\
%    (Fu)^{3} &= 0.
%\end{align}
%%
%Скалярний інваріант~\eqref{eq:I_def} у метриці $(+{-}{-}{-})$:
%%
%\begin{align}\label{eq:I_explicit}
%    \mathcal{I}
%    &= \bigl[(Fu)^{0}\bigr]^2 - \bigl[(Fu)^{1}\bigr]^2 - \bigl[(Fu)^{2}\bigr]^2
%    \notag\\
%    &= \gamma^2\!\left[
%        \frac{E^2 v_x^2}{c^2}
%        - \!\left(E - \frac{Bv_y}{c}\right)^{\!2}
%        - \frac{B^2 v_x^2}{c^2}
%    \right]
%    \notag\\
%    &= \gamma^2\!\left[
%        -E^2\!\left(1 - \frac{v_x^2}{c^2}\right)
%        + \frac{2EBv_y}{c}
%        - \frac{B^2(v_x^2+v_y^2)}{c^2}
%    \right].
%\end{align}
%
%\textbf{5. Подвійна дія тензора поля: $(F^2 u)^{\mu}$.}
%%
%\begin{align}\label{eq:F2u}
%    (F^2 u)^{0} &= F^{0}{}_{1}(Fu)^{1}
%                 = \gamma E\!\left(-E + \frac{Bv_y}{c}\right), \notag\\
%    (F^2 u)^{1} &= F^{1}{}_{0}(Fu)^{0} + F^{1}{}_{2}(Fu)^{2}
%                 = -\frac{\gamma v_x}{c}(E^2 + B^2), \notag\\
%    (F^2 u)^{2} &= F^{2}{}_{1}(Fu)^{1}
%                 = \gamma B\!\left(E - \frac{Bv_y}{c}\right), \notag\\
%    (F^2 u)^{3} &= 0.
%\end{align}
%
%\textbf{6. Просторові компоненти сили ЛЛ.}
%З~\eqref{eq:LL_crossed_simplified} і~\eqref{eq:F2u}, використовуючи
%$u^i = \gamma v^i/c$:
%%
%\begin{align}\label{eq:fLL_components}
%    f^{x}_\text{ЛЛ} &= \frac{2q^4}{3m^2c^5}
%        \left[-\frac{\gamma v_x}{c}(E^2+B^2) - \mathcal{I}\,\frac{\gamma v_x}{c}\right],
%    \notag\\
%    f^{y}_\text{ЛЛ} &= \frac{2q^4}{3m^2c^5}
%        \left[\gamma B\!\left(E - \frac{Bv_y}{c}\right) - \mathcal{I}\,\frac{\gamma v_y}{c}\right].
%\end{align}
%%
%Доданок $\mathcal{I}\,u^i$ в обох компонентах спрямований вздовж $\vect{v}$ ---
%це гальмівна складова сили. Натомість доданок $\gamma B(E - Bv_y/c)\,\hat{y}$
%у другій рядку є перпендикулярним до $\vect{v}$ (при $v_x \neq 0$)
%і відхиляє траєкторію.
%
%\textbf{7. Фізичний зміст.}
%Звичайний $\vect{E}\times\vect{B}$-дрейф:
%%
%\begin{equation}\label{eq:ExB_drift}
%    \vect{v}_d = c\,\frac{\vect{E}\times\vect{B}}{B^2} = \frac{cE}{B}\,\hat{y}.
%\end{equation}
%%
%При $v_y = cE/B$ (режим чистого дрейфу) поперечний доданок
%$\gamma B(E - Bv_y/c)$ у~\eqref{eq:fLL_components} обертається на нуль:
%частинка у режимі дрейфу радіаційного відхилення не відчуває.
%Як тільки $v_y \neq cE/B$ (наявний циклотронний оберт), з'являється
%ненульова поперечна сила, що прагне повернути частинку до дрейфового
%режиму --- ефект, принципово відсутній без урахування реакції випромінювання.
%
%Умова застосовності наближення ЛЛ:
%%
%\begin{equation}\label{eq:LL_validity}
%    \frac{|\vect{f}_\text{ЛЛ}|}{|q\vect{E}|} \sim \tau_0\gamma^2\omega_c \ll 1,
%\end{equation}
%%
%де $\omega_c = qB/(mc)$ --- нерелятивістська циклотронна частота.
%
%\textbf{Відповідь.}
%%
%\begin{equation}\label{eq:fLL_answer}
%    \vect{f}_\text{ЛЛ} = \frac{2q^4\gamma}{3m^2c^6}
%    \left\{
%        -v_x(E^2+B^2)\,\hat{x}
%        + cB\!\left(E - \frac{Bv_y}{c}\right)\hat{y}
%        - \frac{\mathcal{I}\,\vect{v}}{c}
%    \right\},
%\end{equation}
%%
%де $\mathcal{I}$ визначено формулою~\eqref{eq:I_explicit}.
%Перші два доданки відповідають за дрейф траєкторії,
%третій --- за гальмування.
%
\end{problem}